\section{Introduction}
%Elle doit contenir une présentation rapide de l’industrie hydroélectrique, du
%but de faire des essais sur modèle réduit et un descriptif du contenu du rapport.
L’hydroélectricité est l’une des principales sources d’électricité dans le monde, et particulièrement au Québec, où elle représente plus de 80\% de la production électrique \cite{hydroquebec_hydroelectricite}. 
Ces installations reposent sur des turbines dont le type dépend de la hauteur de chute et du débit disponibles. 
La turbine Pelton, une turbine à impulsion, est notamment adaptée aux grandes chutes et aux faibles débits. 
Au Québec et dans le monde entier, l'optimisation de ces turbines représente un enjeu important à l'échelle du parc d'Hydro-Québec, dont la puissance installée peut atteindre jusqu'à 37 200 MW \cite{hydroquebec_hydroelectricite}.
Ainsi, un gain de rendement de seulement 1\% équivaut à plusieurs centaines de mégawatts supplémentaires ce qui devient rapidement important dans un contexte où la demande en électricité ne cesse de croître.
Il est à noter les turbines industrielles sont trop coûteuses et volumineuses pour être optimisées à pleine échelle.
Ainsi, on effectue plutôt des essais surs des modèles réduits et par la suite, on utilise les lois de similitude pour transposer ces résultats au prototype, soit la turbine installée dans une centrale hydroélectrique.
Ce rapport présente d’abord les notions théoriques et le montage expérimental utilisé, puis les résultats obtenus sur une turbine Pelton à échelle réduite.
Ces résultats sont ensuite analysés et discutés avant de conclure sur les principaux résultats.
